
% !TEX program = xelatex
% SDG Equal-Area Sphere Closure Report
% 繁體中文 LaTeX 版本
\documentclass[UTF8,12pt,a4paper]{ctexart}

\usepackage[a4paper,margin=2.4cm]{geometry}
\usepackage{amsmath,amssymb,mathtools,bm}
\usepackage{booktabs,longtable,array}
\usepackage{enumitem}
\usepackage{hyperref}
\usepackage{xcolor}
\usepackage{listings}
\usepackage{caption}

\hypersetup{
  colorlinks=true,
  linkcolor=blue!60!black,
  urlcolor=blue!60!black,
  citecolor=blue!60!black
}

\lstdefinestyle{repo}{
  basicstyle=\ttfamily\small,
  breaklines=true,
  frame=single,
  columns=fullflexible,
  keepspaces=true,
  backgroundcolor=\color{gray!6}
}

\setlist[itemize]{leftmargin=2em}
\setlist[enumerate]{leftmargin=2em}

\title{\textbf{Equal-Area SDG on Sphere 結案報告}\\
\large 極點／Patch Corner 奇異性、等面積投影、守恆性與 Free-Stream Preservation}
\author{專題研究紀錄}
\date{2026-05-05}

\begin{document}
\maketitle

\begin{abstract}
本報告整理本階段針對球面上 Equal-Area SDG 方法所做的設計、實驗、失敗路線排除與最終候選方案。核心問題是原本透過 $A^{-1}$ 進行速度轉換時，在極點或 octahedral patch corner 附近會產生方向不唯一與 numerical divergence 異常。最後採用 stream-function area flux、sphere seam connectivity、edge-streamfunction surface flux、projected surface lifting、global mass correction，以及 Numba fast RHS。此版本已通過質量守恆與 constant-state preservation；短時間 advection convergence 可達高階；但長時間 solid-body rotation 測試約退化到二階上下，推測與 octahedral equal-area map 的 patch corner regularity 與 face-volume flux consistency 有關。
\end{abstract}

\tableofcontents
\newpage

\section{核心結論}

本階段已完成一套可工作的候選方案：
\[
\boxed{
\text{Equal-area map}
+
\text{stream-function area flux}
+
\text{sphere seam connectivity}
+
\text{edge-streamfunction surface flux}
+
\text{global mass correction}
+
\text{Numba fast RHS}
}
\]

最終推薦的 RHS 結構為
\[
R(q)
=
R_{\rm raw}(q)
-
\frac{\int_{S^2} R_{\rm raw}(q)\,dS}
{\int_{S^2}1\,dS},
\]
其中
\[
R_{\rm raw}
=
-\frac{1}{J}
\left[
\partial_x(F^x q)+\partial_y(F^y q)
\right]
+
\frac{1}{J}
M^{-1}
\int_{\partial K}
(f-\hat f)\phi\,ds.
\]

face normal speed 不直接使用兩側各自的 $n\cdot F$，而是由 edge streamfunction trace 給出：
\[
a=F\cdot n=-\partial_s\psi.
\]

這一版已通過：
\begin{enumerate}
  \item \textbf{constant-state preservation}：
  \[
  q=1 \Rightarrow RHS \approx 0.
  \]
  \item \textbf{closed-sphere global mass conservation}：
  \[
  \frac{d}{dt}\int_{S^2}q\,dS\approx0.
  \]
  \item \textbf{sphere seam pairing}：使用 $(X,Y,Z)$ embedding 配對 seam，避免重新引入 $(\lambda,\theta)$ 或 $A^{-1}$ 極點方向不唯一問題。
  \item \textbf{short-time accuracy}：$T=0.1$ Gaussian advection 可觀察到接近四階的 $L^2$ 收斂。
  \item \textbf{工程效能}：Numba fast RHS 與 reference RHS 數值一致，且 RHS benchmark 有約 $40\times$ 至 $120\times$ 加速。
\end{enumerate}

目前尚未完全解決的是
\[
\boxed{
\text{long-time solid-body rotation accuracy 約退化到二階上下。}
}
\]
目前 evidence 顯示這不是 CFL 或 time integration 問題，而是 spatial consistency、seam / patch corner regularity，或 volume-face flux mismatch 長時間累積所造成。

\section{原始問題與研究目標}

\subsection{原始問題}

一開始的主要問題是：
\begin{enumerate}
  \item 需要使用 $A$、$A^{-1}$ 進行球體與八面體 / 平面 patch 間的映射與反映射。
  \item 在極點或 patch corner 附近，反映射可能需要從多個 face 取值，導致極限方向不唯一。
  \item 原本以 $A^{-1}$ 或 velocity transform 得到的速度場，在 $\alpha_0=\pi/4$ 時，divergence diagnostic 不像 $\alpha_0=0$ 那樣漂亮。
  \item 仍然希望保留 \textbf{equal-area mapping}，避免面積權重複雜化或質量積分不自然。
\end{enumerate}

\subsection{本階段目標}

本階段目標可拆成五個：
\begin{enumerate}
  \item 消除 $A^{-1}$ 在極點 / corner 的方向不唯一問題。
  \item 保留 equal-area formulation。
  \item 建立 divergence-free 或 free-stream-preserving 的 velocity / flux 表示。
  \item 建立 closed sphere 上的 conservative DG surface coupling。
  \item 進入可跑 convergence test 的 fast implementation。
\end{enumerate}

\section{第一階段：放棄 $A^{-1}$ Velocity，改用 Stream-Function Area Flux}

\subsection{核心想法}

不再用 $A^{-1}$ 直接把球面速度轉到 flat coordinates，而是使用 stream function：
\[
\psi=-R\,\Omega\cdot X.
\]

在 flat equal-area coordinates 中定義 area flux：
\[
F^x=-\psi_y,\qquad F^y=\psi_x.
\]

這樣
\[
\partial_xF^x+\partial_yF^y
=
-\psi_{yx}+\psi_{xy}
\approx0.
\]

因此 divergence-free property 來自 curl form，而不是來自不穩定的 $A^{-1}$ velocity transform。

\subsection{已驗證結果}

先前 diagnostic 顯示：
\begin{itemize}
  \item \texttt{L2\_stream} 對所有 $\alpha_0$ 都接近 roundoff。
  \item 相較之下，\texttt{L2\_Ainv} 在 $\alpha_0\neq0$ 時約 $10^{-1}$ 等級。
  \item 對 constant volume RHS：
  \[
  RHS\approx0.
  \]
  \item full RHS constant-state check 初步通過。
\end{itemize}

\subsection{判斷}

\[
\boxed{
\text{stream-function area flux 能解決 }\alpha_0=\pi/4\text{ 的 divergence 問題。}
}
\]

\section{第二階段：Projected Surface Lifting}

\subsection{問題}

repo 原本 surface term 有 diagonal lifting 與 projected lifting 兩種可能。為了和 SDG / Table1 quadrature 一致，正式方向應採用 projected inverse mass。

一開始嘗試使用 square Vandermonde 假設：
\[
M^{-1}=VV^T.
\]

但 Table1 order 4, $N=4$ 時：
\[
V.shape=(22,15).
\]

也就是 quadrature / sampling nodes 有 $22$ 個，而 polynomial modes 有 $15$ 個，所以不是 square interpolation system。

\subsection{修正公式}

改用 weighted projection mass：
\[
M=A V^T W V.
\]

projected inverse-mass operator：
\[
P
=
V M^{-1}(A V^T).
\]

實作上避免顯式 inverse：
\[
X=M^{-1}(A V^T),
\qquad
P=VX.
\]

repo lifting convention 是 right multiplication，所以使用
\[
P^T.
\]

\subsection{已驗證結果}

\begin{itemize}
  \item constant field：surface projected lifting 不破壞常數。
  \item \texttt{sphere\_z}：由於跨 element 連續，surface jump 幾乎為 $0$。
  \item \texttt{element\_jump} / \texttt{element\_checker}：surface penalty 被真正啟動，且 projected 與 diagonal lifting 明顯不同。
\end{itemize}

\subsection{判斷}

\[
\boxed{
\text{projected surface lifting 正常。}
}
\]
但後續發現，單純 projected lifting 不是守恆問題的根源。

\section{第三階段：Strong Correction Surface 與 Direct Numerical Flux Surface}

\subsection{初始 Strong Correction 問題}

原本 surface correction 使用
\[
R_{\rm surf}^{strong}
=
\frac{1}{J}
M^{-1}
\int_{\partial K}
(f-\hat f)\phi\,ds.
\]

對 discontinuous field，global mass residual 在 $\alpha_0\neq0$ 時有 $10^{-4}$ 至 $10^{-3}$ 等級。

boundary flux balance diagnostic 顯示：
\begin{itemize}
  \item 非零 mass residual 主要來自 surface term。
  \item 不是 flat open boundary flux 造成。
\end{itemize}

\subsection{Direct Numerical Flux 嘗試}

嘗試改成
\[
R_{\rm surf}^{direct}
=
-\frac{1}{J}
M^{-1}
\int_{\partial K}
\hat f\phi\,ds.
\]

第一版 local-side direct flux 失敗，因為同一條 face 兩側各自算的 $\hat f$ 不滿足
\[
\hat f^-+\hat f^+\approx0.
\]

\texttt{pairErr} 在 $\alpha_0\neq0$ 時約 $10^{-1}$。

\subsection{Pairwise Single-Edge Common Flux}

修正為每條 interior / seam face 只算一次 common flux，再以相反符號分配給左右兩側：
\[
\hat f^-=\hat f,\qquad
\hat f^+=-\hat f.
\]

此後：
\[
\texttt{pairErr}=0.
\]

global conservation balance 通過。

\subsection{Direct Flux 的新問題}

direct numerical flux surface 雖然保住 global mass conservation，但 LSRK constant test 失敗：
\begin{itemize}
  \item 初始 \texttt{rhsM0} 很小。
  \item 但時間推進後 \texttt{massDrift} 與 stage residual 變大。
  \item $q=1$ 被推離常數場。
\end{itemize}

原因是
\[
R_{\rm surf}^{direct}
=
-\frac{1}{J}M^{-1}\int\hat f\phi\,ds
\]
本身不保證 pointwise free-stream preservation。

\subsection{結論}

\[
\boxed{
\text{direct numerical flux surface 不適合作為正式 RHS。}
}
\]
它適合測 conservation，但不適合保持 $q=1$。

\section{第四階段：Sphere Seam Connectivity}

\subsection{核心做法}

flat-square boundary 不再使用 \texttt{same\_state}，而是用球面 embedding $(X,Y,Z)$ 對 boundary faces 進行配對：
\[
\mathbf X^- \approx \mathbf X^+.
\]

這避免使用 $\lambda,\theta$ 或 $A^{-1}$，因此不會引入極點方向不唯一問題。

\subsection{驗證結果}

所有測試中：
\[
\texttt{unmatch}=0,
\]
\[
\texttt{seamErr}\sim10^{-16}.
\]

pairwise flux：
\[
\texttt{pairErr}=0.
\]

closed-sphere mass balance：
\[
I_{\rm full}\approx10^{-14}\sim10^{-12}.
\]

\subsection{判斷}

\[
\boxed{
\text{sphere seam connectivity 成功。}
}
\]

這是整個方案的重要里程碑。

\section{第五階段：Edge-Streamfunction Surface Flux}

\subsection{問題}

direct surface flux 不保 constant state。因此改回 strong correction：
\[
f-\hat f,
\]
但 face normal flux 不再使用 local $n\cdot F$，而是由 edge streamfunction trace 給出：
\[
a=F\cdot n=-\partial_s\psi.
\]

\subsection{RHS 結構}

surface penalty：
\[
p=f-\hat f.
\]

其中：
\[
f=aq.
\]

$\hat f$ 仍採 pairwise common flux 組裝，確保 conservation。

surface RHS：
\[
R_{\rm surf}
=
\frac{1}{J}
M^{-1}
\int_{\partial K}
p\phi\,ds.
\]

\subsection{驗證結果}

此版本通過：
\begin{itemize}
  \item \texttt{constant}：
  \[
  rhs\_L2\sim10^{-14}\sim10^{-12}.
  \]
  \item \texttt{element\_checker} / \texttt{element\_jump}：
  \[
  I_{\rm full}\approx0.
  \]
  \item \texttt{gaussian}：mass residual 接近 roundoff。
  \item edge-pair flux error：
  \[
  \texttt{edgeErr}\approx10^{-15}\sim10^{-13}.
  \]
\end{itemize}

\subsection{判斷}

\[
\boxed{
\text{edge-streamfunction surface flux + strong correction 是目前正確 surface formulation。}
}
\]

\section{第六階段：Local Correction 與 Global Mass Correction}

\subsection{Local Conservative Correction 嘗試}

為了強制每個 element 的 local mass balance，嘗試：
\[
R_K^{corr}
=
R_K^{raw}+c_K.
\]

其中 $c_K$ 使：
\[
\int_K R_K^{corr}\,dS
=
-\int_{\partial K}\hat f\,ds.
\]

\subsection{Local Correction 的問題}

對 $q=1$，local correction 的 \texttt{massDrift} 仍為 roundoff，但 \texttt{corrMax} 在 $\alpha_0\neq0$ 時不是 roundoff。

例如 $nsub=16,\alpha_0=\pi/4$：
\[
\texttt{corrMax}\approx2.4\times10^{-2}.
\]

這代表 local correction 會在 element-wise 層級修改 constant field，存在破壞 free-stream 的風險。

\subsection{結論}

\[
\boxed{
\text{拒絕 local element-wise correction 作為正式版本。}
}
\]

\subsection{Global Mass Correction}

改採全域 correction：
\[
R^{corr}
=
R^{raw}
-
\frac{\int_{S^2}R^{raw}\,dS}
{\int_{S^2}1\,dS}.
\]

這只扣掉全域平均 mass residual，不對每個 element 個別加不同常數。

\subsection{驗證結果}

global correction 版本通過：
\begin{itemize}
  \item constant-state preservation：
  \[
  \texttt{constDev}\sim10^{-14}\sim10^{-13}.
  \]
  \item Gaussian / sphere modes mass drift：
  \[
  \texttt{massDrift}\approx0.
  \]
  \item discontinuous element checker 單步 mass drift：
  \[
  10^{-16}\sim10^{-17}.
  \]
\end{itemize}

\subsection{判斷}

\[
\boxed{
\text{global mass correction 是目前正式候選 correction。}
}
\]

\section{第七階段：Numba Fast RHS}

\subsection{加速目標}

最耗時部分：
\begin{enumerate}
  \item 每次 RHS 的 face edge flux loop。
  \item pairwise penalty assembly。
  \item surface integral accumulation。
  \item LSRK 五 stage 反覆呼叫 RHS。
\end{enumerate}

\subsection{加速策略}

保持數學結構不變，只將 surface-related kernels Numba 化：
\begin{itemize}
  \item edge-streamfunction normal flux kernel；
  \item pairwise strong penalty kernel；
  \item surface integral accumulation kernel；
  \item global mass residual kernel；
  \item global correction application kernel。
\end{itemize}

volume part 仍使用既有 vectorized / matrix operation。

\subsection{Benchmark 結果}

fast RHS 與 reference RHS：
\[
\texttt{diff\_L2}\sim10^{-14}\sim10^{-12}.
\]

mass residual 維持 roundoff。

speedup 約：
\[
40\times\sim120\times.
\]

fast LSRK 也通過：
\begin{itemize}
  \item constant：
  \[
  \texttt{constDev}\sim10^{-14}.
  \]
  \item gaussian：
  \[
  \texttt{massDrift}=0.
  \]
  \item sphere\_z：
  \[
  \texttt{massDrift}\sim10^{-16}.
  \]
  \item element\_checker：
  \[
  \texttt{massDrift}\sim10^{-17}\sim10^{-16}.
  \]
\end{itemize}

\subsection{判斷}

\[
\boxed{
\text{Numba fast RHS 可作為後續 convergence test 的預設實作。}
}
\]

\section{第八階段：Accuracy / Convergence Test}

\subsection{Exact Solution}

solid-body rotation exact solution：
\[
q(X,t)=q_0(R_\Omega(-t)X).
\]

Gaussian initial state：
\[
q_0(X)
=
\exp\left(
-\frac{d(X,X_c)^2}{2\sigma^2}
\right).
\]

\subsection{CFL 控制}

repo 使用 CFL formula：
\[
dt
=
\mathrm{CFL}
\frac{h}{N^2v_{\max}}.
\]

本階段改成 CFL-controlled time step：
\[
nsteps=
\left\lceil
\frac{T}{dt_{\rm nominal}}
\right\rceil,
\qquad
dt=\frac{T}{nsteps}.
\]

\subsection{短時間結果}

對 $T=0.1$、$\alpha_0=\pi/4$、Gaussian：
\[
L^2\text{ rate}\approx4.
\]

代表：
\[
\boxed{
\text{短時間局部 accuracy 是高階的。}
}
\]

\subsection{長時間結果}

對 $T=5.0$：
\begin{itemize}
  \item CFL=1.0 與 CFL=0.5 結果幾乎相同。
  \item 因此 time error 不是主因。
  \item 長時間 $L^2$ rate 約退化到二階上下。
\end{itemize}

\subsection{Gaussian Center 測試}

原始 center：
\[
X_c=(1,0,0)
\]
會靠近 octahedral vertex。

改成：
\[
X_c=\frac{(1,1,1)}{\sqrt3}
\]
後，$T=5$ 仍然約二階上下。

使用其他 generic center 與 sphere modes 時，也觀察到長時間約二階上下。

\subsection{判斷}

\[
\boxed{
\text{長時間 accuracy 退化不是單純 Gaussian center 放在 vertex，也不是 CFL。}
}
\]

目前最可能原因是：
\begin{enumerate}
  \item seam / patch corner interaction；
  \item volume-face flux consistency mismatch；
  \item equal-area octahedral map 在 patch corner 附近的 regularity loss；
  \item edge-streamfunction surface flux 與 volume flux 長時間相位誤差累積。
\end{enumerate}

\section{Face-Flux Consistency Diagnostic}

\subsection{比較對象}

比較
\[
a_{\rm edge}
=
-\partial_s\psi
\]
與
\[
a_{\rm vol}
=
n_xF^x+n_yF^y.
\]

\subsection{結果}

對 $\alpha_0=\pi/4$：
\[
\texttt{diff\_Linf}\approx1.5253\times10^{-1}
\]
幾乎不隨 refinement 下降。

\[
\texttt{diff\_L2}
\sim h^{1/2}.
\]

同時：
\[
\texttt{edgePair}\approx10^{-13},
\]
但
\[
\texttt{volPair}\approx3.05\times10^{-1}.
\]

\subsection{解讀}

\[
\boxed{
a_{\rm edge}\text{ 是 pairwise consistent，但 }a_{\rm vol}\text{ 的 one-sided traces 不形成 common flux。}
}
\]

而 $a_{\rm edge}$ 與 $a_{\rm vol}$ 的 mismatch 有局部 $O(1)$ error，但 measure 很小，所以 weighted $L^2$ 呈 $h^{1/2}$ 下降。

\subsection{Endpoint Diagnostic}

Table1 face trace nodes 不含真正 endpoint，所以 exact endpoint mask 結果為 NaN。top error CSV 顯示最大誤差發生在：
\[
t_{\rm face}\approx0.04691
\quad\text{或}\quad
t_{\rm face}\approx0.95309.
\]

因此最大 mismatch 來自 \textbf{near-corner trace nodes}，不是真正 endpoint。

\subsection{Corner Volavg 嘗試}

嘗試在 near-corner face nodes 以 pairwise averaged volume trace 取代 edge speed：
\[
a_{\rm avg}
=
\frac12(a_{\rm vol}^-+a_{\rm vol}^{+,\text{owner}}).
\]

結果：
\begin{itemize}
  \item $L^\infty$ 略有改善。
  \item $L^2$ 與 convergence rate 沒有改善。
  \item $nsub=16$ 長時間 $L^2$ 仍約 $1.3\times10^{-4}$。
\end{itemize}

\subsection{判斷}

\[
\boxed{
\text{簡單 corner volume-average blending 不是有效修法。}
}
\]

\section{已接受與已拒絕方案}

\subsection{已接受方案}

\begin{longtable}{p{0.27\linewidth}p{0.35\linewidth}p{0.30\linewidth}}
\toprule
模組 & 採用方案 & 理由 \\
\midrule
\endhead
球面到 flat mapping & Equal-area map & 保留面積權重簡潔性 \\
velocity / area flux & Stream-function area flux & 避免 $A^{-1}$ 極點不唯一 \\
seam closure & $(X,Y,Z)$ embedding pairing & 不使用 $\lambda,\theta$，穩定配對 \\
surface lifting & Projected inverse mass & 符合 SDG / Table1 weighted projection \\
face speed & Edge-streamfunction $a=-\partial_s\psi$ & 保守且 free-stream preserving \\
numerical flux & Pairwise common flux & 強制 $\hat f^-+\hat f^+=0$ \\
mass correction & Global correction & 保 mass，不破壞 constant state \\
performance & Numba fast RHS & 數值一致且明顯加速 \\
\bottomrule
\end{longtable}

\subsection{已拒絕方案}

\begin{longtable}{p{0.38\linewidth}p{0.55\linewidth}}
\toprule
方案 & 拒絕原因 \\
\midrule
\endhead
$A^{-1}$ velocity transform & $\alpha_0=\pi/4$ divergence 不穩定，極點方向不唯一 \\
local-side flux assembly & \texttt{pairErr} 在 $\alpha_0\neq0$ 時為 $O(10^{-1})$ \\
direct numerical flux surface & mass 守恆但破壞 constant-state preservation \\
element-wise local correction & 對 constant field 產生非 roundoff correction \\
corner\_volavg face-speed blending & 未改善長時間 $L^2$ rate \\
\bottomrule
\end{longtable}

\section{本階段產生 / 使用的主要檔案}

若尚未 commit，建議將以下檔案整理到 feature branch。

\subsection{Operators}

\begin{lstlisting}[style=repo]
operators/sdg_projected_surface_flux.py
operators/sdg_direct_flux_surface.py
operators/sdg_sphere_seam_connectivity.py
operators/sdg_edge_streamfunction_flux.py
operators/sdg_edge_streamfunction_closed_sphere_rhs.py
operators/sdg_edge_streamfunction_global_corrected_rhs.py
operators/sdg_edge_streamfunction_fast_rhs.py
operators/sdg_hybrid_face_speed.py
\end{lstlisting}

\subsection{Diagnostics / Demos}

\begin{lstlisting}[style=repo]
demo/demo_sdg_projected_surface_lifting_check.py
demo/demo_sdg_boundary_flux_balance_check.py
demo/demo_sdg_direct_flux_surface_balance_check.py
demo/demo_sdg_sphere_seam_flux_balance_check.py
demo/demo_sdg_edge_streamfunction_flux_check.py
demo/demo_sdg_edge_closed_sphere_lsrk_mass_check.py
demo/demo_sdg_edge_global_corrected_lsrk_mass_check.py
demo/demo_sdg_edge_fast_rhs_benchmark.py
demo/demo_sdg_edge_fast_lsrk_mass_check.py
demo/demo_sdg_edge_fast_advection_convergence.py
demo/demo_sdg_edge_fast_advection_convergence_cfl.py
demo/demo_sdg_face_flux_consistency_check.py
demo/demo_sdg_face_flux_endpoint_diagnostic.py
demo/demo_sdg_edge_fast_advection_hybrid_speed_convergence.py
demo/demo_sdg_edge_fast_advection_convergence_cfl_center.py
\end{lstlisting}

\section{目前最終候選演算法}

\subsection{Precompute}

\begin{enumerate}
  \item 建立 equal-area sphere-flat geometry cache。
  \item 建立 SDG Table1 differentiation operators。
  \item 建立 streamfunction：
  \[
  \psi=-R\,\Omega\cdot X.
  \]
  \item 建立 area flux：
  \[
  F^x=-\psi_y,\qquad F^y=\psi_x.
  \]
  \item 建立 projected inverse mass：
  \[
  P=V(A V^TWV)^{-1}(A V^T).
  \]
  \item 用 $(X,Y,Z)$ 建立 sphere seam connectivity。
  \item 建立 face edge derivative matrix。
  \item 預先計算：
  \[
  a_{\rm edge}=-\partial_s\psi.
  \]
\end{enumerate}

\subsection{RHS Evaluation}

給定 $q$：

\begin{enumerate}
  \item Volume:
  \[
  R_{\rm vol}
  =
  -\frac{1}{J}
  \left[
  D_x(F^xq)+D_y(F^yq)
  \right].
  \]

  \item Surface:
  \begin{itemize}
    \item pairwise common numerical flux；
    \item penalty:
    \[
    p=f-\hat f;
    \]
    \item projected lifting:
    \[
    R_{\rm surf}
    =
    \frac{1}{J}
    M^{-1}
    \int_{\partial K}p\phi\,ds.
    \]
  \end{itemize}

  \item Raw:
  \[
  R_{\rm raw}=R_{\rm vol}+R_{\rm surf}.
  \]

  \item Global correction:
  \[
  R
  =
  R_{\rm raw}
  -
  \frac{\int R_{\rm raw}\,dS}{\int1\,dS}.
  \]
\end{enumerate}

\section{尚未解決問題與風險}

\subsection{長時間 Accuracy 退化}

目前最主要限制：
\[
\boxed{
T=5\text{ long-time advection 約二階上下。}
}
\]

這不影響 mass conservation 與 constant-state preservation，但會影響高階 accuracy claim。

\subsection{Volume-Face Flux Mismatch}

diagnostic 顯示
\[
a_{\rm edge}-a_{\rm vol}
\]
在 near-corner trace nodes 有 $O(1)$ mismatch。此 mismatch 可能長時間累積成 phase / amplitude error。

\subsection{Equal-Area Map Regularity}

octahedral equal-area map 在 patch corners 附近可能有限 regularity。即使 solution 本身在 sphere 上 smooth，映射到 flat patch 後可能在 corner / seam 結構下產生低 regularity 表現。

\subsection{Global Correction 對 Accuracy 的影響}

global correction magnitude 通常很小，但仍可能在長時間中對 phase / amplitude 有微弱累積效應。目前未觀察到它是主要問題，但需在報告中標示為限制。

\section{後續建議}

\subsection{短期：整理 Repo}

建議建立 branch：
\begin{lstlisting}[style=repo]
feature/equal-area-edge-streamfunction-sdg
\end{lstlisting}

建議 commit 分組：
\begin{lstlisting}[style=repo]
1. add streamfunction area flux diagnostics
2. add projected surface lifting
3. add sphere seam connectivity
4. add edge-streamfunction surface flux
5. add global mass corrected RHS
6. add numba fast RHS
7. add convergence and face-flux diagnostics
8. add closure report
\end{lstlisting}

\subsection{中期：Accuracy 改善方向}

可嘗試：
\begin{enumerate}
  \item \textbf{face-volume consistency correction}：建立使 volume flux trace 與 edge flux trace 更一致的 split form。
  \item \textbf{corner regularization}：對 near-corner trace nodes 設計更幾何一致的 face speed regularization，而不是簡單 volavg。
  \item \textbf{corner-excluding benchmark}：建立不接近 octahedral vertices 的 manufactured solution / local-in-time benchmark，區分 scheme order 與 mapping singularity effect。
  \item \textbf{multi-patch smoothness analysis}：分析 equal-area map 在 octahedron edge / vertex 的 differentiability class。
  \item \textbf{energy / $L^2$ stability diagnostic}：目前只完整確認 mass 與 constant state，還需測 energy behavior。
\end{enumerate}

\subsection{長期：可寫成專題主軸}

本階段已足以形成一個專題章節：

\begin{quote}
A mass-conservative and free-stream-preserving SDG discretization on an equal-area octahedral sphere map using stream-function based face fluxes.
\end{quote}

核心亮點：
\begin{enumerate}
  \item 避免 $A^{-1}$ pole ambiguity。
  \item 保留 equal-area map。
  \item seam pairing 使用 embedding coordinates。
  \item pairwise flux 保 conservation。
  \item edge-streamfunction face speed 保 free-stream。
  \item Numba fast RHS 可做實驗。
  \item 明確揭示 long-time accuracy 受 patch-corner regularity 限制。
\end{enumerate}

\section{最終摘要}

本對話完成了從「$A^{-1}$ 極點奇異性與 $\alpha_0=\pi/4$ divergence 問題」到「可執行 fast SDG solver prototype」的完整排查與設計。

最終候選版本不是原本的 $A^{-1}$-velocity formulation，也不是 direct flux surface formulation，而是：
\[
\boxed{
\text{equal-area}
+
\text{stream-function volume flux}
+
\text{edge-streamfunction strong surface correction}
+
\text{sphere seam pairing}
+
\text{global mass correction}
+
\text{Numba fast RHS}.
}
\]

此版本已經在 conservation 與 free-stream preservation 上通過測試；短時間 accuracy 高階；長時間 advection 目前觀察到二階上下，推測與 equal-area octahedral patch corner / seam consistency 有關。

因此，本階段可以結案為：
\[
\boxed{
\text{已完成穩定、守恆、可加速的候選方案；下一階段應聚焦 long-time accuracy 與 corner regularity。}
}
\]

\end{document}
